% 第 26 章　温度不是热量（完整章）
\chapter{温度不是热量}

\step{反常识开场}

\begin{mainline}
冬天走进一间没有暖气的教室，伸手摸一下铁制的桌腿，再摸一下木质的桌面。铁桌腿冰凉，木桌面温和，手感差别大到你会脱口而出：“铁比木头冷多了。”可是，如果把一支温度计先后贴在铁桌腿和木桌面上，等读数稳定，你会发现两件东西的温度一模一样——它们在同一间屋子里放了整整一夜，还能差到哪里去？

你的手和温度计，总有一个在说谎。或者更刺激一点：它们都没说谎，只是它们测量的根本不是同一个东西。

再看一件小事。洗澡之前你用手试水温，觉得正好。可如果你刚用冷水洗过手，同一盆水你又会觉得烫；刚用热水洗过手，同一盆水又会觉得凉。水还是那盆水，变的是你的手。我们平时说“我觉得冷、我觉得热”，那个“我”到底在报告什么？

还有更极端的。节日里放的冷焰火、电焊溅出的火花，温度高达上千摄氏度，落在手背上常常只是轻轻一烫，甚至来得及被弹开；而一壶刚烧开的水，温度“只有”一百摄氏度，却没有任何人敢把手伸进去。上千度的东西杀伤力反而不如一百度的——“温度高”和“热能伤人”之间，显然隔着点什么。

这三件事指向同一个结论：我们嘴里的“热”，其实是三个被混为一谈的东西。这一章的任务就是把它们拆开：温度、热量、内能。拆完之后你会发现，日常语言里“这杯水含有很多热量”这种说法，从物理学上看根本不通——不是不准确，是句子本身就有语病。
\end{mainline}

\step{先猜一猜}

照例，先下注再开牌。请把下面三题的答案写下来。

第一题。同一间屋子里放了一夜的一块铁和一块木头，谁的温度低？（a）铁低，因为摸起来凉；（b）木头低，因为木头“存不住热”；（c）一样低，哦不，一样高——一样温度。

第二题。一滴刚烧开的水，和一大壶六十摄氏度的温水，哪一个“含有的热量”更多？（a）开水多，因为它温度高；（b）温水多，因为它量大；（c）这个问题本身问得不对。

第三题。一杯冰水混合物，冰块浮在水面还没化完，放在桌上。假设屋子不算太热，半小时后冰块化掉了一半，这时水的温度是多少？（a）升高了，因为屋子一直在给水加热；（b）还是零摄氏度；（c）升高了一点点，具体多少算不出来。

写下你的答案和理由，本章结尾回来核对。这一章要修理的直觉，是把“温度”“热量”“冷热感觉”当成同一件事的那条直觉。

\step{动手实验}

\begin{experiment}[三杯水和一个说谎的手]
\textbf{材料}：三个一样的杯子、热水（六十摄氏度左右，小心烫）、冷水（加几块冰）、常温水；一支温度计（食品温度计或体温计均可）；一根金属勺、一根木筷、一把塑料尺。

\textbf{步骤一（三杯水）}：左杯装热水，右杯装冷水，中杯装温水（热水冷水各半兑成）。左手放进热水，右手放进冷水，默数六十下。然后两只手同时插进中间的温水。注意听自己的感觉：左手报告“这水是凉的”，右手报告“这水是热的”——两只手泡在同一杯水里，却给出了相反的回答。

\textbf{步骤二（摸三种材料）}：把金属勺、木筷、塑料尺并排放半小时，确保它们真正同温。闭上眼睛，用手背依次快速触摸它们。按“凉感”排序，大多数人会排成：金属最凉、塑料次之、木头最不凉。

\textbf{步骤三（用温度计裁决）}：用温度计分别贴紧三样物品的表面，各等一两分钟。三个读数几乎完全一样，都等于室温。

\textbf{步骤四（做一次预测）}：用温度计测出热水的温度（比如五十五摄氏度）和冰水的温度（比如五摄氏度），各倒半杯，先在心里预测兑匀后的温度，再实际混合测量。大多数人预测三十摄氏度，实测会非常接近——只要动作快、杯子不太厚。把预测值和实测值都记下来，本章的数学部分会告诉你这个“对半取平均”的规则什么时候成立、什么时候会失灵。

\textbf{记录}：写下两个矛盾——左右手对同一杯水的相反判断；手和温度计对同一块金属的不同判断。暂时不用解释，但要记住一个事实：手这个“仪器”，在步骤一里可以被轻易地重新校准，在步骤二里对同温物体给出不同读数。它量的恐怕不是温度。
\end{experiment}

还有几个零成本的日常观察，先存进疑问清单。微波炉里转了两分钟的米饭烫嘴，盛饭的瓷碗却只是温的——同一只炉子里加热，差别从哪里来？发烧时量体温，医生总叮嘱“夹够五分钟”，温度计为什么不能一碰皮肤就出数？冬天对着玻璃窗哈一口气，玻璃上立刻起一层白雾；对镜子哈气，白雾几秒钟又消失——这些“热”的来龙去脉，本章结束时都应该能讲清楚。

\step{旧模型遇到麻烦}

先公平地陈述旧模型。关于“热是什么”，人类历史上最像样的答案是\textbf{热质说}：热是一种看不见、没有质量的流体，叫“热质”。温度高的物体富含热质，温度低的贫乏热质；两个物体一接触，热质就从富处流向贫处，直到两边“水位”拉平。摩擦生热？那是把物体内部藏着的热质“挤”了出来。

这套模型远比它的名声成功。它干净地解释了热接触后冷热中和的现象；在它的框架下，人们发明了量热学，测出了各种物质“容纳热质”的能力，算出的混合温度和实验分毫不差。事实上，直到今天的热学工程计算里，很多公式的形状还留着热质说的影子——它是一条非常、非常好用的错误理论。

但麻烦还是来了，而且是致命的那种。1798 年，伦福德伯爵在慕尼黑兵工厂监督用钝钻头钻炮筒。钻头和炮筒剧烈摩擦，不断变烫，烫到能加热旁边的冷水。按热质说，摩擦只是把金属里存量有限的热质挤出来，可实验里钻头钝到几乎切不下任何碎屑——没有碎屑就没有“被挤出的新表面”——热量却照样源源不断地涌出来，只要马力还在转，热就生个不停。热质不是被挤出来的存货，而是\emph{可以被持续制造}的东西。一种能凭空变出来的“守恒流体”，已经不是流体了。差不多同一时期，戴维做了一个更绝的实验：在抽成真空的容器里，让两块冰互相摩擦，冰居然被磨化了——周围没有任何热源，也没有任何热质可以流进来，热量只能来自摩擦本身。热质说的棺材板被钉上了第二颗钉子。

第二条麻烦来自你自己的手。三杯水实验里，同一杯温水能让左手觉得凉、右手觉得热。如果手是一台测温仪器，那它连“同一输入给出同一输出”都做不到，读数还取决于它刚刚经历过什么。要么承认感觉不可靠，要么承认感觉的本来职责就不是报温度。

第三条麻烦藏在语言里。我们说“热水里含有很多热量”。可搓一搓双手，手就变热了——没有谁给手传过热，热是被\emph{做功}造出来的；再把热手贴在冷墙上，热又流走了。一个“物体里含有”的东西，不该能这样凭空产生、凭空消失。三条麻烦的共同诊断是：热不是一种“东西”，不是一种存货。要走出困境，得先发明一把不受手影响的尺子，再给“热”换一个词性。

\step{发明新概念}

\begin{mainline}
\textbf{新概念一：热平衡。}把两个冷热不同的物体贴在一起，观察它们的宏观性质——体积、电阻、颜色，任何随冷热变化的量。一开始双方都在变，但变化总会停下来，最终进入一种“宏观上再也不变”的状态，这叫\emph{热平衡}。这是一条实验事实：任何材料、任何形状、任何接触方式，只要等得够久，都会到达热平衡。

\textbf{新概念二：热力学第零定律。}现在做一组三个物体的实验：让 A 和温度计 C 达成热平衡，让 B 也和温度计 C 达成热平衡，而 A 和 B 从头到尾没有接触过。这时把 A 和 B 贴到一起——什么都不会发生，它们已经在热平衡中。把这件事写成定律：如果 A 与 C 热平衡，B 与 C 热平衡，那么 A 与 B 热平衡。它听起来像废话，像“等于同量的量彼此相等”的物理版；但它是一条\emph{实验定律}，是无数接触实验的总结，没有一条逻辑命令物体必须这样。它的地位特殊到人们发现：第一、第二定律早已命名，这条更基本的却还没名分，只好插号叫“第零定律”。它真正的用途是发许可证：温度计 C 可以\emph{代表}一切跟它平衡过的物体去互相比较——测温从此可以“隔空传话”，不必让 A 和 B 真的碰面。
\end{mainline}

\begin{center}
\begin{tikzpicture}[scale=1.0]
  % A、B、C 三方与第零定律
  \draw[thick,fill=red!12] (0,0) circle (0.7);
  \node at (0,0) {A};
  \draw[thick,fill=blue!10] (4.4,0) circle (0.7);
  \node at (4.4,0) {B};
  \draw[thick,fill=yellow!25] (2.2,2.2) circle (0.7);
  \node at (2.2,2.2) {C（温度计）};
  \draw[{Stealth}-{Stealth},thick] (0.55,0.62) -- (1.7,1.68);
  \node[above left] at (1.0,1.35) {\small 平衡};
  \draw[{Stealth}-{Stealth},thick] (2.7,1.68) -- (3.85,0.62);
  \node[above right] at (3.4,1.35) {\small 平衡};
  \draw[{Stealth}-{Stealth},thick,dashed] (0.7,0) -- (3.7,0);
  \node[below] at (2.2,-0.75) {\small 则 A 与 B 必然平衡};
\end{tikzpicture}
\end{center}

\begin{mainline}
\textbf{新概念三：温度。}第零定律告诉我们：所有达成热平衡的物体，一定共享某种相同的“东西”。给这个“东西”起个名字，就叫\emph{温度}。注意这个定义的次序：不是先有温度后有平衡，而是先有“平衡”这个可操作的事实，温度只是平衡团体的共同标签。属于同一平衡团体的物体温度相同；不属于的，温度不同——而且实验进一步发现，接触后总是原来“热”的那个变凉、原来“凉”的那个变热，于是温度可以排大小、可以标刻度。温标就是给这个标签配上数字：摄氏温标取冰和水共存的温度为 \(0\,^\circ\mathrm{C}\)、水沸腾为 \(100\,^\circ\mathrm{C}\)，等分一百格；开尔文温标把零点放在能量再也挤不出来的极限处——绝对零度——刻度大小与摄氏度相同，换算就是 \(T = t + 273.15\)。为什么存在一个“低到头”的零点、却没有“高到头”的上限，这个问题会在第 30 章再回来。

有了温度，就轮到温度计出场。温度计的原理都一样：找一种随温度单调变化的性质，把它当“指针”。水银和酒精靠体积胀缩，能读出玻璃管上的刻度；电子体温计靠一种半导体元件电阻随温度的变化；工业上的热电偶靠两种金属接头处产生的微小电压；红外测温枪最特别，它根本不接触被测物，靠接收对方发出的热辐射反推温度——用第 25 章的话说，它在“看”你发出的电磁波。请注意所有温度计的共同点：它们直接显示的永远是\emph{自己}那个敏感元件的温度或状态，是第零定律保证了这个读数可以归到被测物头上。

\textbf{新概念四：微观运动。}温度背后是什么？十九世纪逐步成形的分子动理论给出画面：一切物质都由永不停歇地做无规则运动的微小粒子组成，温度正是这种无规则运动\emph{剧烈程度}的度量——温度越高，微观运动越剧烈。这个画面会在第 27 章被定量地钉死在“分子的平均动能”上，本章只需要它的定性版本：热不是储存在物体里的流体，而是微观粒子“正在进行的运动”的宏观显影。

\textbf{新概念五：内能、热量、功——三个词各管一件事。}这是本章最重要的一次分家。
\begin{itemize}[leftmargin=1.8em]
\item \textbf{内能}：系统内部所有微观粒子的动能和相互作用势能的总和。它是系统的\emph{状态量}——水此刻是多少内能，只取决于它此刻的状态，与它怎么来到这个状态无关。
\item \textbf{热量}：因为温度差而从一个系统转移到另一个系统的能量。它是\emph{过程量}——只发生在转移进行的时候。说“这次加热传了五千焦热量”完全正确；说“这杯水含有五千焦热量”是语病，就像不能说“水库里含有多少下雨”——雨是过程，库里的水是状态。
\item \textbf{功}：改变内能的另一条途径，不靠温度差，靠力学式的推压摩擦——搓手、打气筒压缩空气，都是做功让内能上涨，没有谁给手和打气筒“传热”。
\end{itemize}
这套分家有一个方便的比喻：内能像一个人的存款，热量像别人转给他的转账，功像他加班挣的工资。它像什么：存款是状态，转账和工资是两笔进账的途径，账目上可以分清“余额”和“流水”。它不像什么：钱一旦进账就混在一起了——内能里看不出哪部分是传热来的、哪部分是做功来的，这一点比喻居然也成立；但钱不会像热量那样自发地从富账户流向穷账户，热量流动的方向由温度差单方面决定，这一点比喻就破产了。
\end{mainline}

\step{一句话模型}

\begin{oneline}
温度是状态的标签：处于热平衡的物体温度相同，它度量微观无规则运动的剧烈程度；内能是系统微观能量的总账，是状态量；热量是沿着温度差跨边界转移的那份能量，是过程量——热可以被传递、被制造、被转化，唯独不能被“含有”。说“物体含有热量”，就像说“账户里含有转账”。
\end{oneline}

用这句话检查前面的每个现象：三杯水里，温水是同一个状态、同一个温度，左右手报告不同，说明手报告的从来不是温度；金属和木头同温，手却觉得金属凉，说明手报告的是\emph{能量流出的速度}；火花温度极高但质量极小，内能总账很小，能转移给你的热量也就很小；搓手生热是做功进了内能总账，全程没有热量参与。一句话，全罩住了。

\step{数学翻译器}

直觉到位了，翻译成能算数的语言。核心问题是：传一份热量，温度变多少？

\textbf{第一条式子：比热容。}实验发现，对没有相变的物体，吸收的热量 \(Q\) 与质量 \(m\)、温度变化 \(\Delta T\) 都成正比：
\[
Q \;=\; c\, m\, \Delta T,
\]
比例系数 \(c\) 叫比热容，是材料的性质：一千克这种物质升高一开尔文所需的热量。水的比热容大得出名，\(c_{\text{水}} \approx 4.2\times10^{3}\ \mathrm{J/(kg\cdot K)}\)，是铁（约 \(0.45\times10^{3}\)）的九倍。这条式子在说：同样的热量灌进去，水这种“大肚量”材料只微微升温，铁却明显地热——海边昼夜温差比沙漠小、发动机用水做冷却剂、老式暖水袋灌热水，全是这一个数字在起作用。

公式一写出来，立刻用特殊情况检查。令 \(\Delta T = 0\)：温度不变就不吸放热，合理——没有相变时确实如此。令 \(c\) 极大：同样的 \(Q\) 只换来极小的升温，这正是暖水袋要灌热水而不是热铁块的原因。令 \(Q\) 取负：物体放热、温度下降，式子里的符号自动把方向管起来了。

\textbf{第二条式子：热平衡的账目。}一杯热水和一杯冷水在绝热杯里混合，热水放出的热量等于冷水吸收的热量（这就是当年热质说碰巧正确的领地——没有做功、没有相变时，热量表现得像一笔守恒的账）。设热水质量 \(m_1\)、初温 \(t_1\)，冷水 \(m_2\)、初温 \(t_2\)，混合后温度 \(t\)：
\[
c\,m_1\,(t_1 - t) \;=\; c\,m_2\,(t - t_2).
\]
一千克八十摄氏度的热水兑一千克二十摄氏度的冷水，解出来 \(t = 50\,^\circ\mathrm{C}\)，正好是平均。检查特殊情况：令 \(m_2\) 远大于 \(m_1\)，解出的 \(t\) 几乎等于 \(t_2\)——一滴开水滴进游泳池，池水纹丝不动，合理。

\textbf{第三条式子：相变潜热——温度计罢工的地方。}给冰水混合物持续加热，温度计的读数死死钉在零摄氏度，直到最后一块冰化完，水温才开始爬升。热量一直在流入，温度却不动——按 \(Q=cm\Delta T\) 这是不可能的，说明这条公式的使用前提（没有相变）被打破了。相变期间，热量全部用于拆开冰的微观结构而不是加剧运动，这部分热量叫\emph{潜热}：
\[
Q \;=\; mL,
\]
其中 \(L\) 是单位质量的潜热。冰的熔化热约 \(3.34\times10^{5}\ \mathrm{J/kg}\)，水在一百摄氏度下的汽化热约 \(2.26\times10^{6}\ \mathrm{J/kg}\)——注意单位，这里没有 \(\Delta T\)，因为相变期间温度不变。这条式子解释了一个让直觉出错的事实：一百摄氏度的蒸汽烫伤比一百摄氏度的开水严重得多，因为蒸汽落到皮肤上先液化成水，每克额外甩出两千多焦的潜热，然后才作为一百度的水继续放热。反过来，汽化要\emph{收}潜热，这就是出汗能降温的原因：汗液蒸发时从皮肤抽走热量；也是夏天从泳池出来被风一吹就打哆嗦的原因——风加快了蒸发，潜热的抽水泵开足了马力。

\begin{center}
\begin{tikzpicture}[scale=1.0]
  % 加热曲线：冰→水→汽 的平台图
  \draw[-{Stealth},thick] (0,0) -- (10.2,0) node[below left=2pt and 0pt] {};
  \node[below] at (9.6,-0.05) {\small 吸收的热量 \(Q\)};
  \draw[-{Stealth},thick] (0,0) -- (0,4.6) node[above] {\small 温度};
  % 刻度
  \draw (0,1.15) -- (-0.12,1.15) node[left] {\small \(0\,^\circ\mathrm{C}\)};
  \draw (0,3.35) -- (-0.12,3.35) node[left] {\small \(100\,^\circ\mathrm{C}\)};
  % 曲线
  \draw[very thick,blue!60!black] (0,0.25) -- (1.4,1.15);
  \draw[very thick,blue!60!black] (1.4,1.15) -- (3.2,1.15);
  \draw[very thick,blue!60!black] (3.2,1.15) -- (5.2,3.35);
  \draw[very thick,blue!60!black] (5.2,3.35) -- (8.6,3.35);
  \draw[very thick,blue!60!black,dashed] (8.6,3.35) -- (9.7,4.3);
  % 标注
  \node[below] at (0.7,0.35) {\small 冰升温};
  \node[above] at (2.3,1.15) {\small 冰熔化（温度不变）};
  \node[above,rotate=32] at (4.2,2.1) {\small 水升温};
  \node[above] at (6.9,3.35) {\small 水汽化（温度不变）};
  \node[above,rotate=38] at (9.15,3.95) {\small 蒸汽升温};
\end{tikzpicture}
\end{center}

图里的两段平台就是潜热：热量照收，温度不动。注意汽化的平台远比熔化的长——把一千克一百度的水烧成汽，需要的热量大约是把一千克零度的水烧开的五倍还多。

\textbf{一次带真实数据的估算。}用一只 \(1500\ \mathrm{W}\) 的电热水壶，烧开一升（一千克）二十摄氏度的水，要多久？
\[
Q = c\,m\,\Delta T = 4.2\times10^{3}\times 1\times (100-20) \approx 3.4\times10^{5}\ \mathrm{J}.
\]
功率 \(1500\ \mathrm{W}\) 意味着每秒进账一千五百焦，所以理想时间约 \(3.4\times10^{5}/1.5\times10^{3}\approx 220\) 秒，三四分钟——和厨房里烧水的实际体验吻合（真实的壶会慢一些，因为壶身和空气分走了一部分热量）。同一个数字换个参照系看：三十四万焦大约能把一整千克零摄氏度的冰完全化成零摄氏度的水，或者把十千克的水抬升八十度。能量这笔账，数目字从来不含糊。

\textbf{第四条：热量的三条路。}热量从高温处流向低温处，走路的方式只有三种。
\begin{itemize}[leftmargin=1.8em]
\item \textbf{热传导}：微观粒子互相碰撞接力，把运动的剧烈程度一棒一棒传下去，粒子本身不搬家。金属是接力高手，木头、空气、羽绒是慢性子。大致的规律是：温差越大、接触面积越大、距离越短，传热越快。
\item \textbf{对流}：流体（水、空气）被加热后密度变化、整体流动，连人带马地把能量搬走——粒子亲自搬家。暖气片真正加热房间靠的是它搅动的空气循环，不是它自己那点辐射。
\item \textbf{热辐射}：任何温度的物体都在以电磁波的形式向外发射能量，温度越高发射越强，不需要任何介质，穿越真空畅通无阻——太阳晒热地球走的就是这条路，它和第 25 章的电磁波是同一套物理。
\end{itemize}
一只保温杯就是同时封堵三条路的工程学：双层壁抽真空，掐断传导和对流（没有介质就没有接力和搬家）；内壁镀银，反射辐射；木塞或塑料盖，堵住最后的传导漏口。羽绒服则只对付一件事：用蓬松的绒朵锁住大量静止空气，而静止空气是极好的不良导体——保暖的不是羽绒“产热”，是把你的热量流出的速度压到最低。对流的威力再举两例：海边白天的风几乎总是从海面吹向陆地，因为陆地比热小、升温快，热空气上升后由海面较凉的空气补位，整套环流靠的就是“流体亲自搬家”；烤箱里的热风循环模式让食物熟得更均匀，也是同一个原理。辐射则有一个反直觉的日常版本：深冬夜里站在落地窗前，即便室温不低也觉得“窗那边在吸热”——那不是冷气流进来，而是你的皮肤正隔着玻璃对低温的窗外辐射能量，热量在以光速离开你的身体。

\begin{mathbridge}
辐射的定量规律漂亮得惊人：一块表面单位时间辐射出去的功率为
\[
P \;=\; \varepsilon \sigma A T^{4},
\]
其中 \(T\) 必须是开尔文温度，\(A\) 是表面积，\(\sigma \approx 5.67\times10^{-8}\ \mathrm{W/(m^{2}\cdot K^{4})}\) 是斯忒藩—玻尔兹曼常数，\(\varepsilon\) 是介于 0 和 1 之间的发射率（黑亮粗糙接近 1，银亮抛光接近 0）。那个四次方极其凶猛：温度翻倍，辐射功率变成十六倍。这解释了为什么炼钢炉附近根本无法靠近，也解释了为什么地球平均气温只要波动一点点，辐射出的能量账就会大变。热传导的定量版本同样简洁：单位时间穿过一块材料的热量正比于温差和面积、反比于厚度，比例系数叫热导率——金属的热导率可以是空气的上千倍。顺带一提，任何物体辐射能量的频谱分布只由温度决定，这条“黑体辐射”规律在十九世纪末怎么都对不上理论，正是这场危机逼出了量子论——我们在后面的章节会回到这个事故现场。
\end{mathbridge}

\step{模型的边界}

任何模型都要交代自己的辖区。本章的概念体系有五条重要的边界。

\textbf{边界一：温度是大群体的概念。}温度度量微观运动的剧烈程度，而“剧烈程度”是一个\emph{平均}。对一杯水，平均在亿亿个粒子上，稳如泰山；对十几个粒子的微小系统，平均本身会忽上忽下地涨落，“它的温度是多少”成了一个只有模糊答案的问题；对单个分子，问温度根本没有意义——就像问“这一位观众的平均身高是多少”。温度、压强这些词，天生属于群体。

\textbf{边界二：温度计报告的是它自己的温度。}体温计要夹在腋下几分钟，不是它“反应慢”，而是第零定律需要时间来兑现——温度计必须先和身体达成热平衡，它的读数才等于你的温度。测温枪的读数则是另一套逻辑：它不接触你，靠接收你皮肤发出的红外辐射推算温度，所以它测的其实是“辐射推算值”，隔着玻璃、对着反光金属表面都会出错。

\textbf{边界三：温度可以逐点不同。}火焰上方十厘米和二十厘米，温度差着好几百度；晒太阳的车顶和车底，温度也完全不同。“温度”必须逐点定义，只有达成热平衡的区域内部，才谈得上一个共同的温度。太空中一块被太阳照着的卫星，向阳面和背阴面可以相差上百摄氏度，说“这颗卫星的温度”必须指明问的是哪里。

\textbf{边界四：比热容并非常数。}\(Q=cm\Delta T\) 里的 \(c\) 只在有限的温度范围内近似不变。温度降到极低时，各种材料的比热都会急剧变小，趋向于零——这背后是纯量子力学的原因，和“边界五”一起指向同一个事实：本章的比热模型是常温世界的近似，不是终极规律。

\textbf{边界五：绝对零度只许逼近，不许可达。}开尔文温标的零点不是一个普通的冷点，它是微观运动能抽取的能量的下限。实验和理论共同表明：任何制冷手段都只能一步步逼近绝对零度，不能真正到达。这条禁令的完整理由需要热力学第三定律，第 30 章会把它讲透。

\begin{challenge}
温度的现代定义比“热的程度”更深刻：在统计物理里，温度定义为内能随熵变化的变化率的倒数。这个定义有一个匪夷所思的合法推论——在某些特殊系统里（能量有上限的系统，比如外磁场中的一堆原子核自旋），当粒子被“泵”到高能态的数目超过低能态时，温度可以取\emph{负值}，而且负温度比任何正温度都“更热”：负温度系统与正温度系统接触，热量从负温度流向正温度。激光器里的“粒子数反转”正是这种状态的工作基础。这个例子提醒一件事：日常直觉里“温度就是冷热的刻度”只是温度概念的起点，不是它的终点。
\end{challenge}

\step{回头解释开场现象}

现在验收新模型，逐个清算开场的三个悬念。

\textbf{铁桌腿为什么比木桌面凉？}它们温度相同——温度计没有说谎，第零定律保证一夜之后整间屋子都在热平衡里。说谎的是“手=温度计”这条默认假设。手的皮肤里有专门的感觉神经，它们报告的不是皮肤此刻的温度，而是\emph{热量流出皮肤的速度}。手摸铁桌腿时，铁的热导率高，热量被迅速从接触点抽走、扩散到整块桌腿，皮肤局部温度骤降，神经大喊“凉”；摸木头时，木头是热的不良导体，热量抽走一点就堵在表面，皮肤降温慢，于是“温和”。同温不同感，差别全在传热速度上。所以手从来就不是温度计，它是一台简陋的\emph{热流计}——这个发现不丢人，第 28 章会看到，连精密的量热学也是从“跟踪热流”起家的。

\textbf{三杯水为什么自相矛盾？}同理。手在热水里泡了一分钟，皮肤升温、向外散热达到一个平衡；插进温水时，温水比皮肤凉，热量从皮肤流向水，这台热流计报告“凉”。另一只手刚从冷水出来，皮肤温度低于温水，热量流入皮肤，报告“热”。两台热流计被预先校准到了相反的状态，对同一杯水自然给出相反读数。第零定律的意义正在于此：它给了我们不依赖任何“手感”的仲裁方案——温度计不需要先泡手。

\textbf{火花为什么不伤人？}温度、内能、热量三个概念各司其职的最漂亮演示。火花的温度确实高达上千摄氏度——微观运动极其剧烈；但每颗火花只是极小的一粒金属或氧化物，质量也许只有万分之一克量级，内能总账极小。它落到皮肤上，能转移过来的热量微乎其微，只够烫热一个针尖大的点。一壶开水温度“只有”一百度，可一千克水的内能总账巨大，倒在皮肤上就是几万焦的热量灌进去。杀伤力看的是转移到你身上的\emph{热量}，不看火花的\emph{温度}——这正是标题那句话的含义：温度不是热量。

至于“先猜一猜”的三题：第一题，同温，手感和温度是两回事；第二题，问题本身就是语病，正确的问题是“哪个内能大”（壶大得多）和“哪个与皮肤接触会传递更多热量”（也是壶）；第三题，只要冰没化完，冰水混合物被相变钉在零摄氏度——屋子传进来的热量全被潜热吃掉了，温度纹丝不动。

顺手清算实验后存下的三笔账。体温计要等几分钟，是因为第零定律不是瞬时生效的——温度计得先和你达成热平衡，它报的本来就是自己的温度。微波炉里饭烫碗温，是因为微波主要被水分子吸收（第 25 章讲过电磁波如何与物质相互作用），米饭含水多、瓷碗含水少，两者吸收的能量天差地别；之后碗里那点温度，有一部分还是从米饭经传导“借”来的。哈气起雾则是潜热的现身说法：呼出的气体里满是水汽，遇到冷玻璃液化成细小液滴——你看到的白雾不是气体，是液体小珠；几秒后小液滴重新蒸发，从玻璃上抽走汽化热，雾就散了。

\step{脑洞实验与挑战题}

\begin{thought}[把一支温度计扔到深空]
想象把一支足够结实的温度计带出太阳系，扔到远离任何恒星的星际空间，然后等——等足够久。它最终会显示多少度？没有空气接触它，传导和对流都不存在，但辐射这条路永远开着：温度计不断向外辐射能量降温，同时也接收来自四面八方的辐射。最终它会和“整个宇宙的辐射场”达成热平衡，停在约 \(2.7\ \mathrm{K}\)——宇宙大爆炸遗留下来的微波背景辐射的温度。这个思想实验说明两件事：第一，辐射让“热平衡”可以跨越真空发生，第零定律不需要接触；第二，温度的标尺长得离谱——实验室里的冷原子已经被压到十亿分之一开尔文以下，而大质量恒星核心的温度在一千万开尔文以上，人类迄今在加速器里撞出的瞬间高温超过万亿开尔文。从最低到最高，温度这条轴横跨二十多个数量级，而每一格上住着的物质形态都完全不同。
\end{thought}

\begin{puzzles}
\item 一块八十摄氏度的铁块和一块二十摄氏度、等质量的铁块接触，平衡温度是五十摄氏度。现在把冷铁块换成等质量、二十摄氏度的\emph{水}，平衡温度会高于五十还是低于五十？（提示：比较铁和水的比热容，想一想“谁更难被温度打动”。答案是远低于五十——实际约三十摄氏度上下。）
\item 在青藏高原上烧水，水早早就沸腾翻滚了，面条却总煮不透。是火不够旺吗？加大火力有用吗？（提示：沸腾是被相变钉住的温度平台，气压低时这个平台出现在更低的温度；火再旺，水温也不会超过沸点，多余的热量全变成蒸汽带走了。高压锅的用途正是人为抬高这个平台。）
\item 夏天烈日下暴晒一小时的铁栏杆和木栏杆，摸上去的差别比冬天更大还是更小？（提示：把手当成热流计重新分析，注意这次热流的方向反了过来——再想想，暴晒下的金属和木头真的同温吗？辐射和传导一起决定了答案比冬天更复杂。）
\item 把一块零摄氏度的冰扔进一杯零摄氏度的水里，盖上绝热盖子，冰会融化吗？如果把整杯放到零摄氏度的房间里呢？（提示：先检查有没有温度差——没有温差就没有热量流动；再想想潜热从哪里来。）
\end{puzzles}

这一章我们拆开了温度、热量、内能三个词，但只立了账本的格式，还没有立账本的总规则：传热和做功这两笔进账，加起来如何决定内能的变化？为什么有的账能做、有的账永远做不成？那是第 28 章和第 29 章要立的两条定律。
